\chapter{Background}

This chapter introduces the two problems addressed by the thesis ---
automatic news summarization and automatic fact-checking --- the
Mixture-of-Agents (MoA) technique that supports the technical contribution,
and a survey of related work in Vietnamese and international NLP. The chapter
ends with a statement of the gap this thesis closes.

\section{Automatic Text Summarization}

\textbf{Definition.}
Automatic text summarization is the task of producing a shorter version of a
source text that preserves its essential information. Given an input document
$D$, a summarizer $f$ produces a summary $S = f(D)$ such that $|S| \ll |D|$
and $S$ retains the key facts of $D$.

\textbf{Classification.}
Summarization systems are commonly classified along two axes.

\begin{itemize}
  \item \textit{Extractive vs.\ abstractive.}
        Extractive systems select existing sentences (or sub-sentences) from
        the source. Abstractive systems generate new text and can paraphrase,
        re-order, or combine ideas. Modern LLM-based summarizers are
        predominantly abstractive.
  \item \textit{Single-document vs.\ multi-document.}
        Single-document systems summarize one article in isolation;
        multi-document systems aggregate information across several related
        articles. The system in this thesis is single-document, but the MoA
        pipeline can be read as a degenerate multi-source aggregator in
        which the ``multiple sources'' are different LLM drafts of the same
        article.
\end{itemize}

\textbf{Vietnamese-specific challenges.}
Vietnamese poses three difficulties for summarization.
\begin{enumerate}
  \item \textit{Word segmentation.} Vietnamese is written with a space
        between syllables, not between words; a single conceptual word may
        span two or three syllables (\textit{tóm tắt}, \textit{kiểm chứng}).
        Token-overlap metrics computed on whitespace tokens systematically
        under-credit paraphrases that recompose multi-syllable terms.
  \item \textit{Diacritics and tone marks.} Vietnamese has six lexical tones,
        encoded as diacritics on vowels; encoder vocabularies that fragment
        diacriticised characters lose semantic granularity.
  \item \textit{Scarcity of evaluation corpora.} Compared with English
        (CNN/DailyMail, XSum, MultiNews), Vietnamese has no widely-adopted,
        human-annotated summarization benchmark with reference summaries.
        This is the practical reason ROUGE/BLEU/BERTScore in Vietnamese
        research are usually computed \textit{against the source article},
        not against human reference summaries --- the methodological caveat
        revisited in Chapter~4.
\end{enumerate}

\section{News Fact-Checking}

\textbf{Misinformation in Vietnamese media.}
Vietnamese readers encounter clickbait headlines, partial quotations, and
outright misinformation across social media and news aggregators. Manual
fact-checking does not scale; an automated complement that can flag
unsupported claims is increasingly valuable.

\textbf{The standard pipeline.}
The literature converges on a three-step automatic fact-checking pipeline.
\begin{enumerate}
  \item \textit{Claim detection.} Extract check-worthy atomic claims from a
        document.
  \item \textit{Evidence retrieval.} For each claim, retrieve passages from
        a corpus (or the open web) that may support or refute it.
  \item \textit{Verdict prediction.} Classify each claim as \textit{entailed},
        \textit{contradicted}, or \textit{not mentioned} given the evidence,
        and aggregate per-claim verdicts into a per-document score.
\end{enumerate}

\textbf{Search-augmented verification.}
Open-web search has become the faithful evidence source for
fact-checking, because no curated knowledge base keeps pace with breaking
news. Search-augmented verification combines a web-search API with an LLM
that reads the retrieved snippets and renders a verdict; this is the design
adopted in Chapter~3 (Tavily $\rightarrow$ GPT-4o-mini).

\section{Mixture-of-Agents (MoA)}

\textbf{Origin.}
Wang et al.\ proposed Mixture-of-Agents in 2024 \cite{Wang2024}. The central
empirical observation is that an LLM acting as an \textit{aggregator} produces
better answers when given several candidate answers from other LLMs ---
even when those other LLMs are individually weaker than the aggregator
itself. The authors call this the \textit{collaborativeness of LLMs}.

\textbf{Architecture.}
MoA is a layered structure. Layer $\ell$ consists of $n$ agents
$\{A_{\ell,1}, \ldots, A_{\ell,n}\}$; each agent at layer $\ell+1$ receives
the prompt together with the outputs of all agents at layer $\ell$. Formally,
agent $i$ at layer $\ell+1$ produces

\[
y_{\ell+1,i} \;=\; A_{\ell+1,i}\!\left(\,[\,y_{\ell,1}, y_{\ell,2}, \ldots, y_{\ell,n}\,] \,+\, x \right),
\]

\noindent where $x$ is the user prompt and $+$ denotes concatenation. In a
two-layer instantiation (the one used in this thesis), layer~1 is the
\textit{proposer} layer and layer~2 is a single \textit{aggregator}.

\textbf{Reported results.}
On the AlpacaEval~2.0 benchmark, an MoA configuration with six open-source
proposers and one aggregator reaches an LC win rate of $65.1\%$, surpassing
GPT-4o ($57.5\%$) under the same evaluation protocol. Performance scales
monotonically with the number of layers and the number of proposers per
layer (Tables~3--4 of \cite{Wang2024}).

\textbf{Open question for this thesis.}
The MoA paper evaluates on instruction-following benchmarks (AlpacaEval,
MT-Bench, FLASK) where the LLM judge --- not overlap against any source
text --- is the ground truth. Whether MoA transfers to \textit{grounded}
summarization (where a source article exists and overlap can
be measured), and whether the benefit shows up in Vietnamese, is the
empirical question this thesis answers.

\section{Related Works}

\textbf{Vietnamese summarization systems.}
\textit{Google News Vietnam} and \textit{Báo Mới} provide aggregated news with
short snippets, but their snippets are extractive lead-paragraph excerpts,
not editorial abstractive summaries. Academic Vietnamese summarization work
has used ViT5 \cite{Nguyen2022} for abstractive
generation, evaluated almost exclusively with ROUGE against either reference
summaries (when available) or the source article (more often).

\textbf{Fact-checking tools.}
\textit{ClaimBuster} (UT Arlington) detects check-worthy claims in English
political speech; the \textit{Google Fact Check Tools} API surfaces existing
fact-checking-organization verdicts. Neither targets Vietnamese, and neither
operates at the article level for ordinary readers.

\textbf{Browser-extension reading aids.}
\textit{TLDR This}, \textit{Summari}, and \textit{Reader Mode} extensions
inject in-page summaries on top of arbitrary web pages but are
English-first; their summarizers degrade on Vietnamese, and they do not
offer fact-checking. \textit{NewsGuard} and \textit{Bias Buddy} score the
reliability of \textit{publishers}, not individual articles, and rely on
manual ratings rather than per-article verification.

\textbf{LLM-as-judge evaluation.}
Zheng et al.\ \cite{Zheng2023} introduced MT-Bench and the Chatbot Arena
methodology, demonstrating that a strong LLM (GPT-4) agrees with human
preferences on chatbot outputs at $\sim$80\% (broadly the same level at
which two humans agree with each other). FLASK \cite{Ye2023} extended this
to fine-grained rubric scoring along multiple alignment skills.
AlpacaEval \cite{Dubois2024} added position randomization and length
control to mitigate two known LLM-judge biases. The Axis~B evaluation
module of this thesis is built on these methods.

\textbf{The gap.}
To the best of my knowledge, no prior work has (i)~applied Mixture-of-Agents
to Vietnamese, (ii)~combined extension-based product delivery with
fact-checking and summarization on Vietnamese news, or (iii)~explicitly
contrasted overlap, LLM-judge, and human evaluation in a single
\textit{three-axis} framework for Vietnamese summarization. This thesis fills
all three gaps.

\section{Summary}

This chapter framed Vietnamese news summarization and fact-checking as the
two parallel needs the thesis addresses, introduced Mixture-of-Agents as the
technique that the thesis adapts, surveyed related Vietnamese and
international work, and identified the gap: no prior Vietnamese system
combines MoA fusion, browser-extension delivery, and a methodological three-axis evaluation. Chapter~2 develops the theoretical
foundations that the system in Chapter~3 builds upon.
